% =====================================================================
%  Thème 2 — Chasles et combinaisons — Niveau FACILE — Exercices
% =====================================================================
\documentclass[11pt,a4paper]{article}
\usepackage[utf8]{inputenc}
\usepackage[T1]{fontenc}
\usepackage{lmodern}
\usepackage[french]{babel}
\usepackage{amsmath,amssymb,mathtools}
\usepackage{xcolor}
\usepackage{enumitem}
\usepackage{fancyhdr}
\usepackage[a4paper,margin=2.1cm,headheight=26pt,headsep=10pt]{geometry}

\definecolor{primary}{RGB}{222,70,80}
\definecolor{primarydark}{RGB}{150,30,42}

\pagestyle{fancy}\fancyhf{}
\fancyhead[L]{\small\textcolor{primarydark}{\textbf{Fiche 2} — Calcul vectoriel}}
\fancyhead[R]{\small\textcolor{primarydark}{\textsc{Thème 2 — Facile — Exercices}}}
\fancyfoot[C]{\small\thepage}
\renewcommand{\headrulewidth}{0.8pt}

\newcommand{\vc}[1]{\overrightarrow{#1}}
\providecommand{\norm}[1]{\left\lVert #1\right\rVert}
\newcommand{\exo}[1]{\medskip\noindent{\color{primary}\bfseries\large Exercice #1.}\par\nopagebreak\smallskip}

\begin{document}
\begin{center}
{\LARGE\bfseries\color{primarydark}Chasles et combinaisons — Niveau Facile}\\[2pt]
{\color{primary}\rule{0.5\linewidth}{1.1pt}}
\end{center}
\vspace{1mm}
{\small\itshape Objectif : appliquer $\vc{AB}+\vc{BC}=\vc{AC}$, reconnaître le vecteur nul et les
opposés. On ne calcule pas de coordonnées : on manipule les lettres.}
\vspace{2mm}

\exo{1} \textbf{Compléter par la relation de Chasles.}
Simplifier chaque somme en un seul vecteur.
\begin{enumerate}[label=\textbf{\alph*)},leftmargin=2em,itemsep=2pt]
  \item $\vc{AB}+\vc{BC}$
  \item $\vc{MN}+\vc{NP}$
  \item $\vc{RS}+\vc{ST}+\vc{TU}$
  \item $\vc{EF}+\vc{FG}+\vc{GH}+\vc{HK}$
\end{enumerate}

\exo{2} \textbf{Vecteur nul et opposés.}
Donner le résultat le plus simple de chaque expression.
\begin{enumerate}[label=\textbf{\alph*)},leftmargin=2em,itemsep=2pt]
  \item $\vc{AB}+\vc{BA}$
  \item $\vc{AB}+\vc{BC}+\vc{CA}$
  \item $-\vc{DE}$ (l'exprimer comme un vecteur à lettres inversées)
  \item $\vc{PQ}+\vc{QP}+\vc{PQ}$
\end{enumerate}

\exo{3} \textbf{Transformer une soustraction.}
Simplifier en utilisant $-\vc{XY}=\vc{YX}$.
\begin{enumerate}[label=\textbf{\alph*)},leftmargin=2em,itemsep=2pt]
  \item $\vc{AB}-\vc{CB}$
  \item $\vc{AC}-\vc{AB}$
  \item $\vc{MP}-\vc{MN}$
\end{enumerate}

\exo{4} \textbf{Choisir un point intermédiaire.}
Dans chaque cas, décomposer le vecteur en passant par le point indiqué.
\begin{enumerate}[label=\textbf{\alph*)},leftmargin=2em,itemsep=2pt]
  \item Écrire $\vc{AC}$ en passant par $B$.
  \item Écrire $\vc{AD}$ en passant par $B$ puis $C$.
  \item Écrire $\vc{MN}$ en passant par un point $O$ quelconque.
\end{enumerate}

\exo{5} \textbf{Parallélogramme.}
$ABCD$ est un parallélogramme (sommets dans l'ordre). Compléter les égalités.
\begin{enumerate}[label=\textbf{\alph*)},leftmargin=2em,itemsep=2pt]
  \item $\vc{AB}=\vc{D\;?}$
  \item $\vc{AD}=\vc{B\;?}$
  \item $\vc{AB}+\vc{AD}=\vc{A\;?}$ (règle du parallélogramme)
\end{enumerate}

\end{document}
